% Volume VII: Formal Synthesis
% Part XIII. A Mathematics of Disciplined Distrust
% Chapter 73: Evidentiary Reach
% Spine: proof-spines/ch073-spine.md

\chapter{Evidentiary Reach}
\label{ch:ch073}

This chapter formalizes \textbf{evidentiary reach} as the bounded scope within which a datum materially changes warranted belief. For evidence \(E\), the reach set is
\[
\mathcal R(E)
=
\left\{
H:
\inf_{M\in\mathfrak M}
\left|
P_M(H\mid E)-P_M(H)
\right|>\epsilon
\right\}.
\]
\Definition\ A proposition \(H\) belongs to \(\mathcal R(E)\) only if, across the class \(\mathfrak M\) of robustly admissible models, conditioning on \(E\) changes its likelihood by more than \(\epsilon\). Reach is therefore not vividness, nor mere rhetorical force; it is tested inferential movement that survives model variation.

The distinction matters because evidence can strongly bear on one claim while failing to justify transfers into broader narratives, character judgments, or policy conclusions. This chapter thus gives formal content to the rule that the magnitude of an update should track tested evidentiary reach rather than the emotional intensity of a record. Once \(\mathcal R(E)\) is defined, the next chapter can bound not only where evidence reaches, but how much motion it may authorize.
